% =====================================================================
% Epoch 07 — Pillars 4-7 advancement + first Stage-1 scorecard snapshot
%   (epoch retrospective; then-current status, NOT current canonical row)
% Status: [DRAFT] [NEEDS_UPDATE] Wave 7 epoch retrospective
% Window: 2026-04-22 to 2026-04-23
% Math notes: Math70-84 range
% AUDIT-FLAG: AUDIT-2026-05-02-Wave7-Aux-Epoch-Overclaim
%   Hostile-referee audit (2026-05-02) flagged abstract phrasing
%   ("completed a critical synthesis", "Pillar 4 gauge interactions
%   and SO(10) uniqueness", "Pillar 5 chirality closure",
%   "Pillar 6 generational structure", "Pillar 7 quantum consistency")
%   as compressing the canonical-tier story into "closed". Math302
%   formalized the carve-out separating proof-rigor and
%   realisation-status; Math270-313 defence and verdict cycles came
%   AFTER this epoch. Title and abstract recast as a then-current
%   advancement and first stage-wide scorecard snapshot, NOT a
%   current definitive status.
% =====================================================================
% --- TECT canonical preamble (see ../../papers/_shared/tect-preamble.tex)
% =====================================================================
% TECT canonical LaTeX preamble (revtex4-2 / single-column / theorem-safe)
% =====================================================================
% Maintainer: Jusang Lee (jtkor@outlook.com)
% Binding from: 2026-05-06
% Purpose: a single source-of-truth preamble for every TECT paper
% (Paper-00 .. Paper-10 + Auxiliary notes). Designed to (a) eliminate
% font-corruption on pdflatex (T1 fontenc + lmodern), (b) make
% theorem-heavy mathematical-physics content render cleanly in a
% single-column layout (PRD class option, NOT PRL two-column), (c)
% provide consistent theorem numbering across all TECT papers via
% amsthm with a shared counter set, (d) make hyperref + cleveref
% cross-references resolve cleanly with the standard two-pass
% pdflatex compile.
%
% Usage in each Paper-NN.tex:
%   % =====================================================================
% TECT canonical LaTeX preamble (revtex4-2 / single-column / theorem-safe)
% =====================================================================
% Maintainer: Jusang Lee (jtkor@outlook.com)
% Binding from: 2026-05-06
% Purpose: a single source-of-truth preamble for every TECT paper
% (Paper-00 .. Paper-10 + Auxiliary notes). Designed to (a) eliminate
% font-corruption on pdflatex (T1 fontenc + lmodern), (b) make
% theorem-heavy mathematical-physics content render cleanly in a
% single-column layout (PRD class option, NOT PRL two-column), (c)
% provide consistent theorem numbering across all TECT papers via
% amsthm with a shared counter set, (d) make hyperref + cleveref
% cross-references resolve cleanly with the standard two-pass
% pdflatex compile.
%
% Usage in each Paper-NN.tex:
%   % =====================================================================
% TECT canonical LaTeX preamble (revtex4-2 / single-column / theorem-safe)
% =====================================================================
% Maintainer: Jusang Lee (jtkor@outlook.com)
% Binding from: 2026-05-06
% Purpose: a single source-of-truth preamble for every TECT paper
% (Paper-00 .. Paper-10 + Auxiliary notes). Designed to (a) eliminate
% font-corruption on pdflatex (T1 fontenc + lmodern), (b) make
% theorem-heavy mathematical-physics content render cleanly in a
% single-column layout (PRD class option, NOT PRL two-column), (c)
% provide consistent theorem numbering across all TECT papers via
% amsthm with a shared counter set, (d) make hyperref + cleveref
% cross-references resolve cleanly with the standard two-pass
% pdflatex compile.
%
% Usage in each Paper-NN.tex:
%   \input{../_shared/tect-preamble.tex}
%   \begin{document}
%   ...
%   \end{document}
%
% Build (every paper): `pdflatex Paper-NN && pdflatex Paper-NN`
% (the second pass resolves \ref{} / \cref{} forward references).
% A bash/PowerShell wrapper that automates this is at
% Docs/papers/papers/_shared/build-paper.sh / build-paper.ps1.
% =====================================================================

% ---- Document class --------------------------------------------------
% PRD class option of revtex4-2 with [reprint,onecolumn] gives the same
% APS heading / by-line / abstract style as PRL but in a wider single
% column that handles long display equations and theorem environments
% without column-break artefacts. nofootinbib keeps citations out of
% footnotes. The 11pt option restores standard font metrics; with T1
% fontenc + lmodern this gives non-bitmap glyphs.
\documentclass[%
  aps,prd,reprint,onecolumn,nofootinbib,11pt,%
  amsmath,amssymb,floatfix,longbibliography%
]{revtex4-2}

% ---- Encoding & fonts ------------------------------------------------
\usepackage[T1]{fontenc}        % proper 8-bit font encoding (no font corruption)
\usepackage[utf8]{inputenc}     % allow UTF-8 source
\usepackage{lmodern}            % scalable Latin Modern (replaces bitmap CM)
\usepackage{textcomp}           % extra text symbols
\usepackage{microtype}          % subtle kerning improvements

% ---- UTF-8 → LaTeX character mappings --------------------------------
% Maps the Unicode characters that occur in TECT body text (legacy from
% the chat-content auto-archival rule, CLAUDE.md §4) to LaTeX-safe
% commands. Without these declarations, T1+inputenc rejects the chars
% with "Unicode character X not set up for use with LaTeX".
% Adding declarations centrally (here) is preferable to per-file edits.
\DeclareUnicodeCharacter{00A7}{\S}              % §  (section sign)
\DeclareUnicodeCharacter{00B2}{\ensuremath{^{2}}}        % ²
\DeclareUnicodeCharacter{00B3}{\ensuremath{^{3}}}        % ³
\DeclareUnicodeCharacter{00B5}{\ensuremath{\mu}}         % µ (micro)
\DeclareUnicodeCharacter{00B7}{\ensuremath{\cdot}}       % · (middle dot)
\DeclareUnicodeCharacter{00D7}{\ensuremath{\times}}      % ×
\DeclareUnicodeCharacter{00E4}{\"a}             % ä
\DeclareUnicodeCharacter{00E9}{\'e}             % é
\DeclareUnicodeCharacter{00F6}{\"o}             % ö
\DeclareUnicodeCharacter{00FC}{\"u}             % ü
\DeclareUnicodeCharacter{010C}{\v{C}}           % Č
\DeclareUnicodeCharacter{010D}{\v{c}}           % č
\DeclareUnicodeCharacter{0394}{\ensuremath{\Delta}}      % Δ
\DeclareUnicodeCharacter{03A3}{\ensuremath{\Sigma}}      % Σ
\DeclareUnicodeCharacter{03A9}{\ensuremath{\Omega}}      % Ω
\DeclareUnicodeCharacter{03B1}{\ensuremath{\alpha}}      % α
\DeclareUnicodeCharacter{03B2}{\ensuremath{\beta}}       % β
\DeclareUnicodeCharacter{03B3}{\ensuremath{\gamma}}      % γ
\DeclareUnicodeCharacter{03B4}{\ensuremath{\delta}}      % δ
\DeclareUnicodeCharacter{03B5}{\ensuremath{\varepsilon}} % ε
\DeclareUnicodeCharacter{03BA}{\ensuremath{\kappa}}      % κ
\DeclareUnicodeCharacter{03BB}{\ensuremath{\lambda}}     % λ
\DeclareUnicodeCharacter{03BC}{\ensuremath{\mu}}         % μ
\DeclareUnicodeCharacter{03BD}{\ensuremath{\nu}}         % ν
\DeclareUnicodeCharacter{03C0}{\ensuremath{\pi}}         % π
\DeclareUnicodeCharacter{03C1}{\ensuremath{\rho}}        % ρ
\DeclareUnicodeCharacter{03C3}{\ensuremath{\sigma}}      % σ
\DeclareUnicodeCharacter{03C4}{\ensuremath{\tau}}        % τ
\DeclareUnicodeCharacter{03C6}{\ensuremath{\varphi}}     % φ
\DeclareUnicodeCharacter{03C7}{\ensuremath{\chi}}        % χ
\DeclareUnicodeCharacter{03C8}{\ensuremath{\psi}}        % ψ
\DeclareUnicodeCharacter{03C9}{\ensuremath{\omega}}      % ω
\DeclareUnicodeCharacter{2013}{--}              % – (en dash)
\DeclareUnicodeCharacter{2014}{---}             % — (em dash)
\DeclareUnicodeCharacter{2018}{`}               % ‘
\DeclareUnicodeCharacter{2019}{'}               % ’
\DeclareUnicodeCharacter{201C}{``}              % “
\DeclareUnicodeCharacter{201D}{''}              % ”
\DeclareUnicodeCharacter{2070}{\ensuremath{^{0}}}        % ⁰
\DeclareUnicodeCharacter{2071}{\ensuremath{^{i}}}        % ⁱ
\DeclareUnicodeCharacter{2122}{\textsuperscript{TM}}     % ™
\DeclareUnicodeCharacter{2126}{\ensuremath{\Omega}}      % Ω (ohm sign)
\DeclareUnicodeCharacter{210F}{\ensuremath{\hbar}}       % ℏ
\DeclareUnicodeCharacter{2115}{\ensuremath{\mathbb{N}}}  % ℕ
\DeclareUnicodeCharacter{2102}{\ensuremath{\mathbb{C}}}  % ℂ
\DeclareUnicodeCharacter{211D}{\ensuremath{\mathbb{R}}}  % ℝ
\DeclareUnicodeCharacter{2124}{\ensuremath{\mathbb{Z}}}  % ℤ
\DeclareUnicodeCharacter{211A}{\ensuremath{\mathbb{Q}}}  % ℚ
\DeclareUnicodeCharacter{2192}{\ensuremath{\to}}         % →
\DeclareUnicodeCharacter{21D2}{\ensuremath{\Rightarrow}} % ⇒
\DeclareUnicodeCharacter{2200}{\ensuremath{\forall}}     % ∀
\DeclareUnicodeCharacter{2203}{\ensuremath{\exists}}     % ∃
\DeclareUnicodeCharacter{2208}{\ensuremath{\in}}         % ∈
\DeclareUnicodeCharacter{2209}{\ensuremath{\notin}}      % ∉
\DeclareUnicodeCharacter{2227}{\ensuremath{\wedge}}      % ∧
\DeclareUnicodeCharacter{2228}{\ensuremath{\vee}}        % ∨
\DeclareUnicodeCharacter{2229}{\ensuremath{\cap}}        % ∩
\DeclareUnicodeCharacter{222A}{\ensuremath{\cup}}        % ∪
\DeclareUnicodeCharacter{2245}{\ensuremath{\cong}}       % ≅
\DeclareUnicodeCharacter{2248}{\ensuremath{\approx}}     % ≈
\DeclareUnicodeCharacter{2260}{\ensuremath{\neq}}        % ≠
\DeclareUnicodeCharacter{2264}{\ensuremath{\leq}}        % ≤
\DeclareUnicodeCharacter{2265}{\ensuremath{\geq}}        % ≥
\DeclareUnicodeCharacter{22C5}{\ensuremath{\cdot}}       % ⋅
\DeclareUnicodeCharacter{2295}{\ensuremath{\oplus}}      % ⊕
\DeclareUnicodeCharacter{2297}{\ensuremath{\otimes}}     % ⊗

% ---- Mathematics & symbols ------------------------------------------
\usepackage{amsmath,amssymb,bm,mathrsfs,mathtools}
\usepackage{amsthm}             % theorem environments (load BEFORE hyperref)

% ---- Theorem environments (shared counter set, TECT canonical) -------
% All theorems / lemmas / propositions / corollaries / remarks share a
% single counter so cross-references read consistently across a paper.
% Definitions get a separate counter (definitions are numbered in their
% own sequence by editorial convention).
\newtheorem{theorem}{Theorem}
\newtheorem{lemma}[theorem]{Lemma}
\newtheorem{proposition}[theorem]{Proposition}
\newtheorem{corollary}[theorem]{Corollary}

\theoremstyle{remark}
\newtheorem{remark}[theorem]{Remark}
\newtheorem{example}[theorem]{Example}

\theoremstyle{definition}
\newtheorem{definition}[theorem]{Definition}
\newtheorem{conjecture}[theorem]{Conjecture}
\newtheorem{hypothesis}[theorem]{Hypothesis}

% ---- Tables / floats / graphics --------------------------------------
\usepackage{booktabs}           % \toprule / \midrule / \bottomrule
\usepackage{array}
\usepackage{graphicx}
\usepackage{enumitem}           % \begin{itemize}[leftmargin=...] options
\usepackage{hyphenat}           % \hyp{} for breakable hyphens in \texttt
\usepackage{seqsplit}           % \seqsplit{} for long unbreakable strings
\sloppy                         % allow stretchier inter-word spacing to
                                % reduce overfull \hbox warnings on long URLs
                                % and long \texttt tokens (paragraph-level).
\emergencystretch=3em           % last-resort hbox stretch budget

% ---- Cross-references (hyperref last among the standard packages,
%      cleveref AFTER hyperref) ----------------------------------------
\usepackage[%
  colorlinks=true,%
  linkcolor=blue,%
  citecolor=blue,%
  urlcolor=blue,%
  breaklinks=true,%
  bookmarks=true,%
  bookmarksnumbered=true%
]{hyperref}
\usepackage[capitalise,nameinlink]{cleveref}

% ---- TECT-canonical author block macros -----------------------------
% (defined here so every paper has the same author / affiliation style)
\newcommand{\tectauthor}{%
  \author{Jusang Lee}%
  \email{jtkor@outlook.com}%
  \affiliation{Independent researcher, Republic of Korea}%
}

% ---- TECT-canonical archive footer ----------------------------------
\newcommand{\tectarchivefooter}{%
  \par\medskip
  \noindent\small\textit{Canonical archive}:
  \url{https://github.com/TECT-OpenPhysics/TECT}.\par%
}

% ---- TECT TODO / AUDIT-FLAG / HONEST-SCOPE inline marks --------------
% Used in drafts to highlight pending audit items without disturbing
% the body text style. Suppress via \tectsuppressmarks (define empty).
\providecommand{\tectauditflag}[1]{\textcolor{red}{\textbf{[AUDIT-FLAG: #1]}}}
\providecommand{\tecttodo}[1]{\textcolor{orange}{\textbf{[TODO: #1]}}}
\providecommand{\tecthonestscope}[1]{\textit{Honest scope: #1.}}

% =====================================================================
% End of TECT canonical preamble
% =====================================================================

%   \begin{document}
%   ...
%   \end{document}
%
% Build (every paper): `pdflatex Paper-NN && pdflatex Paper-NN`
% (the second pass resolves \ref{} / \cref{} forward references).
% A bash/PowerShell wrapper that automates this is at
% Docs/papers/papers/_shared/build-paper.sh / build-paper.ps1.
% =====================================================================

% ---- Document class --------------------------------------------------
% PRD class option of revtex4-2 with [reprint,onecolumn] gives the same
% APS heading / by-line / abstract style as PRL but in a wider single
% column that handles long display equations and theorem environments
% without column-break artefacts. nofootinbib keeps citations out of
% footnotes. The 11pt option restores standard font metrics; with T1
% fontenc + lmodern this gives non-bitmap glyphs.
\documentclass[%
  aps,prd,reprint,onecolumn,nofootinbib,11pt,%
  amsmath,amssymb,floatfix,longbibliography%
]{revtex4-2}

% ---- Encoding & fonts ------------------------------------------------
\usepackage[T1]{fontenc}        % proper 8-bit font encoding (no font corruption)
\usepackage[utf8]{inputenc}     % allow UTF-8 source
\usepackage{lmodern}            % scalable Latin Modern (replaces bitmap CM)
\usepackage{textcomp}           % extra text symbols
\usepackage{microtype}          % subtle kerning improvements

% ---- UTF-8 → LaTeX character mappings --------------------------------
% Maps the Unicode characters that occur in TECT body text (legacy from
% the chat-content auto-archival rule, CLAUDE.md §4) to LaTeX-safe
% commands. Without these declarations, T1+inputenc rejects the chars
% with "Unicode character X not set up for use with LaTeX".
% Adding declarations centrally (here) is preferable to per-file edits.
\DeclareUnicodeCharacter{00A7}{\S}              % §  (section sign)
\DeclareUnicodeCharacter{00B2}{\ensuremath{^{2}}}        % ²
\DeclareUnicodeCharacter{00B3}{\ensuremath{^{3}}}        % ³
\DeclareUnicodeCharacter{00B5}{\ensuremath{\mu}}         % µ (micro)
\DeclareUnicodeCharacter{00B7}{\ensuremath{\cdot}}       % · (middle dot)
\DeclareUnicodeCharacter{00D7}{\ensuremath{\times}}      % ×
\DeclareUnicodeCharacter{00E4}{\"a}             % ä
\DeclareUnicodeCharacter{00E9}{\'e}             % é
\DeclareUnicodeCharacter{00F6}{\"o}             % ö
\DeclareUnicodeCharacter{00FC}{\"u}             % ü
\DeclareUnicodeCharacter{010C}{\v{C}}           % Č
\DeclareUnicodeCharacter{010D}{\v{c}}           % č
\DeclareUnicodeCharacter{0394}{\ensuremath{\Delta}}      % Δ
\DeclareUnicodeCharacter{03A3}{\ensuremath{\Sigma}}      % Σ
\DeclareUnicodeCharacter{03A9}{\ensuremath{\Omega}}      % Ω
\DeclareUnicodeCharacter{03B1}{\ensuremath{\alpha}}      % α
\DeclareUnicodeCharacter{03B2}{\ensuremath{\beta}}       % β
\DeclareUnicodeCharacter{03B3}{\ensuremath{\gamma}}      % γ
\DeclareUnicodeCharacter{03B4}{\ensuremath{\delta}}      % δ
\DeclareUnicodeCharacter{03B5}{\ensuremath{\varepsilon}} % ε
\DeclareUnicodeCharacter{03BA}{\ensuremath{\kappa}}      % κ
\DeclareUnicodeCharacter{03BB}{\ensuremath{\lambda}}     % λ
\DeclareUnicodeCharacter{03BC}{\ensuremath{\mu}}         % μ
\DeclareUnicodeCharacter{03BD}{\ensuremath{\nu}}         % ν
\DeclareUnicodeCharacter{03C0}{\ensuremath{\pi}}         % π
\DeclareUnicodeCharacter{03C1}{\ensuremath{\rho}}        % ρ
\DeclareUnicodeCharacter{03C3}{\ensuremath{\sigma}}      % σ
\DeclareUnicodeCharacter{03C4}{\ensuremath{\tau}}        % τ
\DeclareUnicodeCharacter{03C6}{\ensuremath{\varphi}}     % φ
\DeclareUnicodeCharacter{03C7}{\ensuremath{\chi}}        % χ
\DeclareUnicodeCharacter{03C8}{\ensuremath{\psi}}        % ψ
\DeclareUnicodeCharacter{03C9}{\ensuremath{\omega}}      % ω
\DeclareUnicodeCharacter{2013}{--}              % – (en dash)
\DeclareUnicodeCharacter{2014}{---}             % — (em dash)
\DeclareUnicodeCharacter{2018}{`}               % ‘
\DeclareUnicodeCharacter{2019}{'}               % ’
\DeclareUnicodeCharacter{201C}{``}              % “
\DeclareUnicodeCharacter{201D}{''}              % ”
\DeclareUnicodeCharacter{2070}{\ensuremath{^{0}}}        % ⁰
\DeclareUnicodeCharacter{2071}{\ensuremath{^{i}}}        % ⁱ
\DeclareUnicodeCharacter{2122}{\textsuperscript{TM}}     % ™
\DeclareUnicodeCharacter{2126}{\ensuremath{\Omega}}      % Ω (ohm sign)
\DeclareUnicodeCharacter{210F}{\ensuremath{\hbar}}       % ℏ
\DeclareUnicodeCharacter{2115}{\ensuremath{\mathbb{N}}}  % ℕ
\DeclareUnicodeCharacter{2102}{\ensuremath{\mathbb{C}}}  % ℂ
\DeclareUnicodeCharacter{211D}{\ensuremath{\mathbb{R}}}  % ℝ
\DeclareUnicodeCharacter{2124}{\ensuremath{\mathbb{Z}}}  % ℤ
\DeclareUnicodeCharacter{211A}{\ensuremath{\mathbb{Q}}}  % ℚ
\DeclareUnicodeCharacter{2192}{\ensuremath{\to}}         % →
\DeclareUnicodeCharacter{21D2}{\ensuremath{\Rightarrow}} % ⇒
\DeclareUnicodeCharacter{2200}{\ensuremath{\forall}}     % ∀
\DeclareUnicodeCharacter{2203}{\ensuremath{\exists}}     % ∃
\DeclareUnicodeCharacter{2208}{\ensuremath{\in}}         % ∈
\DeclareUnicodeCharacter{2209}{\ensuremath{\notin}}      % ∉
\DeclareUnicodeCharacter{2227}{\ensuremath{\wedge}}      % ∧
\DeclareUnicodeCharacter{2228}{\ensuremath{\vee}}        % ∨
\DeclareUnicodeCharacter{2229}{\ensuremath{\cap}}        % ∩
\DeclareUnicodeCharacter{222A}{\ensuremath{\cup}}        % ∪
\DeclareUnicodeCharacter{2245}{\ensuremath{\cong}}       % ≅
\DeclareUnicodeCharacter{2248}{\ensuremath{\approx}}     % ≈
\DeclareUnicodeCharacter{2260}{\ensuremath{\neq}}        % ≠
\DeclareUnicodeCharacter{2264}{\ensuremath{\leq}}        % ≤
\DeclareUnicodeCharacter{2265}{\ensuremath{\geq}}        % ≥
\DeclareUnicodeCharacter{22C5}{\ensuremath{\cdot}}       % ⋅
\DeclareUnicodeCharacter{2295}{\ensuremath{\oplus}}      % ⊕
\DeclareUnicodeCharacter{2297}{\ensuremath{\otimes}}     % ⊗

% ---- Mathematics & symbols ------------------------------------------
\usepackage{amsmath,amssymb,bm,mathrsfs,mathtools}
\usepackage{amsthm}             % theorem environments (load BEFORE hyperref)

% ---- Theorem environments (shared counter set, TECT canonical) -------
% All theorems / lemmas / propositions / corollaries / remarks share a
% single counter so cross-references read consistently across a paper.
% Definitions get a separate counter (definitions are numbered in their
% own sequence by editorial convention).
\newtheorem{theorem}{Theorem}
\newtheorem{lemma}[theorem]{Lemma}
\newtheorem{proposition}[theorem]{Proposition}
\newtheorem{corollary}[theorem]{Corollary}

\theoremstyle{remark}
\newtheorem{remark}[theorem]{Remark}
\newtheorem{example}[theorem]{Example}

\theoremstyle{definition}
\newtheorem{definition}[theorem]{Definition}
\newtheorem{conjecture}[theorem]{Conjecture}
\newtheorem{hypothesis}[theorem]{Hypothesis}

% ---- Tables / floats / graphics --------------------------------------
\usepackage{booktabs}           % \toprule / \midrule / \bottomrule
\usepackage{array}
\usepackage{graphicx}
\usepackage{enumitem}           % \begin{itemize}[leftmargin=...] options
\usepackage{hyphenat}           % \hyp{} for breakable hyphens in \texttt
\usepackage{seqsplit}           % \seqsplit{} for long unbreakable strings
\sloppy                         % allow stretchier inter-word spacing to
                                % reduce overfull \hbox warnings on long URLs
                                % and long \texttt tokens (paragraph-level).
\emergencystretch=3em           % last-resort hbox stretch budget

% ---- Cross-references (hyperref last among the standard packages,
%      cleveref AFTER hyperref) ----------------------------------------
\usepackage[%
  colorlinks=true,%
  linkcolor=blue,%
  citecolor=blue,%
  urlcolor=blue,%
  breaklinks=true,%
  bookmarks=true,%
  bookmarksnumbered=true%
]{hyperref}
\usepackage[capitalise,nameinlink]{cleveref}

% ---- TECT-canonical author block macros -----------------------------
% (defined here so every paper has the same author / affiliation style)
\newcommand{\tectauthor}{%
  \author{Jusang Lee}%
  \email{jtkor@outlook.com}%
  \affiliation{Independent researcher, Republic of Korea}%
}

% ---- TECT-canonical archive footer ----------------------------------
\newcommand{\tectarchivefooter}{%
  \par\medskip
  \noindent\small\textit{Canonical archive}:
  \url{https://github.com/TECT-OpenPhysics/TECT}.\par%
}

% ---- TECT TODO / AUDIT-FLAG / HONEST-SCOPE inline marks --------------
% Used in drafts to highlight pending audit items without disturbing
% the body text style. Suppress via \tectsuppressmarks (define empty).
\providecommand{\tectauditflag}[1]{\textcolor{red}{\textbf{[AUDIT-FLAG: #1]}}}
\providecommand{\tecttodo}[1]{\textcolor{orange}{\textbf{[TODO: #1]}}}
\providecommand{\tecthonestscope}[1]{\textit{Honest scope: #1.}}

% =====================================================================
% End of TECT canonical preamble
% =====================================================================

%   \begin{document}
%   ...
%   \end{document}
%
% Build (every paper): `pdflatex Paper-NN && pdflatex Paper-NN`
% (the second pass resolves \ref{} / \cref{} forward references).
% A bash/PowerShell wrapper that automates this is at
% Docs/papers/papers/_shared/build-paper.sh / build-paper.ps1.
% =====================================================================

% ---- Document class --------------------------------------------------
% PRD class option of revtex4-2 with [reprint,onecolumn] gives the same
% APS heading / by-line / abstract style as PRL but in a wider single
% column that handles long display equations and theorem environments
% without column-break artefacts. nofootinbib keeps citations out of
% footnotes. The 11pt option restores standard font metrics; with T1
% fontenc + lmodern this gives non-bitmap glyphs.
\documentclass[%
  aps,prd,reprint,onecolumn,nofootinbib,11pt,%
  amsmath,amssymb,floatfix,longbibliography%
]{revtex4-2}

% ---- Encoding & fonts ------------------------------------------------
\usepackage[T1]{fontenc}        % proper 8-bit font encoding (no font corruption)
\usepackage[utf8]{inputenc}     % allow UTF-8 source
\usepackage{lmodern}            % scalable Latin Modern (replaces bitmap CM)
\usepackage{textcomp}           % extra text symbols
\usepackage{microtype}          % subtle kerning improvements

% ---- UTF-8 → LaTeX character mappings --------------------------------
% Maps the Unicode characters that occur in TECT body text (legacy from
% the chat-content auto-archival rule, CLAUDE.md §4) to LaTeX-safe
% commands. Without these declarations, T1+inputenc rejects the chars
% with "Unicode character X not set up for use with LaTeX".
% Adding declarations centrally (here) is preferable to per-file edits.
\DeclareUnicodeCharacter{00A7}{\S}              % §  (section sign)
\DeclareUnicodeCharacter{00B2}{\ensuremath{^{2}}}        % ²
\DeclareUnicodeCharacter{00B3}{\ensuremath{^{3}}}        % ³
\DeclareUnicodeCharacter{00B5}{\ensuremath{\mu}}         % µ (micro)
\DeclareUnicodeCharacter{00B7}{\ensuremath{\cdot}}       % · (middle dot)
\DeclareUnicodeCharacter{00D7}{\ensuremath{\times}}      % ×
\DeclareUnicodeCharacter{00E4}{\"a}             % ä
\DeclareUnicodeCharacter{00E9}{\'e}             % é
\DeclareUnicodeCharacter{00F6}{\"o}             % ö
\DeclareUnicodeCharacter{00FC}{\"u}             % ü
\DeclareUnicodeCharacter{010C}{\v{C}}           % Č
\DeclareUnicodeCharacter{010D}{\v{c}}           % č
\DeclareUnicodeCharacter{0394}{\ensuremath{\Delta}}      % Δ
\DeclareUnicodeCharacter{03A3}{\ensuremath{\Sigma}}      % Σ
\DeclareUnicodeCharacter{03A9}{\ensuremath{\Omega}}      % Ω
\DeclareUnicodeCharacter{03B1}{\ensuremath{\alpha}}      % α
\DeclareUnicodeCharacter{03B2}{\ensuremath{\beta}}       % β
\DeclareUnicodeCharacter{03B3}{\ensuremath{\gamma}}      % γ
\DeclareUnicodeCharacter{03B4}{\ensuremath{\delta}}      % δ
\DeclareUnicodeCharacter{03B5}{\ensuremath{\varepsilon}} % ε
\DeclareUnicodeCharacter{03BA}{\ensuremath{\kappa}}      % κ
\DeclareUnicodeCharacter{03BB}{\ensuremath{\lambda}}     % λ
\DeclareUnicodeCharacter{03BC}{\ensuremath{\mu}}         % μ
\DeclareUnicodeCharacter{03BD}{\ensuremath{\nu}}         % ν
\DeclareUnicodeCharacter{03C0}{\ensuremath{\pi}}         % π
\DeclareUnicodeCharacter{03C1}{\ensuremath{\rho}}        % ρ
\DeclareUnicodeCharacter{03C3}{\ensuremath{\sigma}}      % σ
\DeclareUnicodeCharacter{03C4}{\ensuremath{\tau}}        % τ
\DeclareUnicodeCharacter{03C6}{\ensuremath{\varphi}}     % φ
\DeclareUnicodeCharacter{03C7}{\ensuremath{\chi}}        % χ
\DeclareUnicodeCharacter{03C8}{\ensuremath{\psi}}        % ψ
\DeclareUnicodeCharacter{03C9}{\ensuremath{\omega}}      % ω
\DeclareUnicodeCharacter{2013}{--}              % – (en dash)
\DeclareUnicodeCharacter{2014}{---}             % — (em dash)
\DeclareUnicodeCharacter{2018}{`}               % ‘
\DeclareUnicodeCharacter{2019}{'}               % ’
\DeclareUnicodeCharacter{201C}{``}              % “
\DeclareUnicodeCharacter{201D}{''}              % ”
\DeclareUnicodeCharacter{2070}{\ensuremath{^{0}}}        % ⁰
\DeclareUnicodeCharacter{2071}{\ensuremath{^{i}}}        % ⁱ
\DeclareUnicodeCharacter{2122}{\textsuperscript{TM}}     % ™
\DeclareUnicodeCharacter{2126}{\ensuremath{\Omega}}      % Ω (ohm sign)
\DeclareUnicodeCharacter{210F}{\ensuremath{\hbar}}       % ℏ
\DeclareUnicodeCharacter{2115}{\ensuremath{\mathbb{N}}}  % ℕ
\DeclareUnicodeCharacter{2102}{\ensuremath{\mathbb{C}}}  % ℂ
\DeclareUnicodeCharacter{211D}{\ensuremath{\mathbb{R}}}  % ℝ
\DeclareUnicodeCharacter{2124}{\ensuremath{\mathbb{Z}}}  % ℤ
\DeclareUnicodeCharacter{211A}{\ensuremath{\mathbb{Q}}}  % ℚ
\DeclareUnicodeCharacter{2192}{\ensuremath{\to}}         % →
\DeclareUnicodeCharacter{21D2}{\ensuremath{\Rightarrow}} % ⇒
\DeclareUnicodeCharacter{2200}{\ensuremath{\forall}}     % ∀
\DeclareUnicodeCharacter{2203}{\ensuremath{\exists}}     % ∃
\DeclareUnicodeCharacter{2208}{\ensuremath{\in}}         % ∈
\DeclareUnicodeCharacter{2209}{\ensuremath{\notin}}      % ∉
\DeclareUnicodeCharacter{2227}{\ensuremath{\wedge}}      % ∧
\DeclareUnicodeCharacter{2228}{\ensuremath{\vee}}        % ∨
\DeclareUnicodeCharacter{2229}{\ensuremath{\cap}}        % ∩
\DeclareUnicodeCharacter{222A}{\ensuremath{\cup}}        % ∪
\DeclareUnicodeCharacter{2245}{\ensuremath{\cong}}       % ≅
\DeclareUnicodeCharacter{2248}{\ensuremath{\approx}}     % ≈
\DeclareUnicodeCharacter{2260}{\ensuremath{\neq}}        % ≠
\DeclareUnicodeCharacter{2264}{\ensuremath{\leq}}        % ≤
\DeclareUnicodeCharacter{2265}{\ensuremath{\geq}}        % ≥
\DeclareUnicodeCharacter{22C5}{\ensuremath{\cdot}}       % ⋅
\DeclareUnicodeCharacter{2295}{\ensuremath{\oplus}}      % ⊕
\DeclareUnicodeCharacter{2297}{\ensuremath{\otimes}}     % ⊗

% ---- Mathematics & symbols ------------------------------------------
\usepackage{amsmath,amssymb,bm,mathrsfs,mathtools}
\usepackage{amsthm}             % theorem environments (load BEFORE hyperref)

% ---- Theorem environments (shared counter set, TECT canonical) -------
% All theorems / lemmas / propositions / corollaries / remarks share a
% single counter so cross-references read consistently across a paper.
% Definitions get a separate counter (definitions are numbered in their
% own sequence by editorial convention).
\newtheorem{theorem}{Theorem}
\newtheorem{lemma}[theorem]{Lemma}
\newtheorem{proposition}[theorem]{Proposition}
\newtheorem{corollary}[theorem]{Corollary}

\theoremstyle{remark}
\newtheorem{remark}[theorem]{Remark}
\newtheorem{example}[theorem]{Example}

\theoremstyle{definition}
\newtheorem{definition}[theorem]{Definition}
\newtheorem{conjecture}[theorem]{Conjecture}
\newtheorem{hypothesis}[theorem]{Hypothesis}

% ---- Tables / floats / graphics --------------------------------------
\usepackage{booktabs}           % \toprule / \midrule / \bottomrule
\usepackage{array}
\usepackage{graphicx}
\usepackage{enumitem}           % \begin{itemize}[leftmargin=...] options
\usepackage{hyphenat}           % \hyp{} for breakable hyphens in \texttt
\usepackage{seqsplit}           % \seqsplit{} for long unbreakable strings
\sloppy                         % allow stretchier inter-word spacing to
                                % reduce overfull \hbox warnings on long URLs
                                % and long \texttt tokens (paragraph-level).
\emergencystretch=3em           % last-resort hbox stretch budget

% ---- Cross-references (hyperref last among the standard packages,
%      cleveref AFTER hyperref) ----------------------------------------
\usepackage[%
  colorlinks=true,%
  linkcolor=blue,%
  citecolor=blue,%
  urlcolor=blue,%
  breaklinks=true,%
  bookmarks=true,%
  bookmarksnumbered=true%
]{hyperref}
\usepackage[capitalise,nameinlink]{cleveref}

% ---- TECT-canonical author block macros -----------------------------
% (defined here so every paper has the same author / affiliation style)
\newcommand{\tectauthor}{%
  \author{Jusang Lee}%
  \email{jtkor@outlook.com}%
  \affiliation{Independent researcher, Republic of Korea}%
}

% ---- TECT-canonical archive footer ----------------------------------
\newcommand{\tectarchivefooter}{%
  \par\medskip
  \noindent\small\textit{Canonical archive}:
  \url{https://github.com/TECT-OpenPhysics/TECT}.\par%
}

% ---- TECT TODO / AUDIT-FLAG / HONEST-SCOPE inline marks --------------
% Used in drafts to highlight pending audit items without disturbing
% the body text style. Suppress via \tectsuppressmarks (define empty).
\providecommand{\tectauditflag}[1]{\textcolor{red}{\textbf{[AUDIT-FLAG: #1]}}}
\providecommand{\tecttodo}[1]{\textcolor{orange}{\textbf{[TODO: #1]}}}
\providecommand{\tecthonestscope}[1]{\textit{Honest scope: #1.}}

% =====================================================================
% End of TECT canonical preamble
% =====================================================================


\begin{document}

\title{TECT epoch~7 retrospective (Math70--84, 2026-04-22 to 04-23):
Pillars 4--7 advancement and first stage-wide scorecard snapshot}

\author{Jusang Lee}
\email{jtkor@outlook.com}
\affiliation{Independent researcher, Republic of Korea}

\date{\today}

\begin{abstract}
This paper is an EPOCH RETROSPECTIVE compiling the internal theorem
programme established by the Math70--84 archive of the Topological
Energy Condensate Theory (TECT) during 2026-04-22 to 04-23. This
epoch consolidates the then-current advancement of Pillars 4--7 and
records the first stage-wide scorecard, without claiming that the
later Math270--313 defence and verdict cycles were already complete.
The four advancement themes recorded in the Math70--84 archive were:
candidate Pillar~4 gauge-group emergence via defect-bundle topology
on $\mathbb{CP}^2$ and equivariant-moduli arguments; candidate
Pillar~5 chirality consolidation via index-theoretic protection;
candidate Pillar~6 generational structure via a Pati--Salam RGE +
Yukawa hierarchy; and candidate Pillar~7 quantum consistency via
Ward-identity and modular-arithmetic checks. The first stage-wide
tier-coherent snapshot was a then-current 4 PROVED (T7) / 2 PROVED
CONDITIONAL (T6) / 1 CLOSED@1-LOOP (T5) / 3 STRONG EVIDENCE (T4) /
1 OPEN-NEGATIVE (T0+T2) summary; the current canonical scorecard is
recorded in \texttt{Docs/status/TOE-FACT-SHEET.md} and reflects the
Math270--313 defence cycles and the Math302 proof-rigor /
realisation-gate carve-out. This document does NOT advance the
canonical tier.
\end{abstract}

\maketitle

% =====================================================================
\section{Introduction}
% =====================================================================
The TECT framework rests on an 11-pillar foundation (Math54), each addressing a
distinct sector of fundamental physics: emergence of the mass gap, Lorentz
invariance, spin-2 gravity, electroweak gauge symmetry, chirality, generations,
quantum consistency, transversality, equivalence principle, $\hbar$-origin, and
the cosmological constant. After four Epochs of foundational work (Math01-69),
the research programme required a systematic tier assignment reflecting the
current mathematical status of each pillar. Epoch 7 delivers that assessment and
simultaneously advances four pillars to new closure depths.

This paper documents the chronological arc of 2026-04-22 to 04-23, chronicling
(i) the Pillar 4 gauge-group programme under Math75 + Math80 (defect moduli,
equivariance, SO(10) uniqueness, Yukawa hierarchy), (ii) the Pillar 5 chirality
closure via Math76 (index-theoretic protection), (iii) the Pillar 6 generational
embedding via Math77 (Pati-Salam + RGE), (iv) the Pillar 7 quantum consistency
via Ward + modular constraints, and (v) the integration of all 11 pillars into a
canonical Stage-1 scorecard (Math78 synthesis, later audit-rolled-back to
PARTIAL-ADVANCED but directionally sound).

% =====================================================================
\section{Setting and Pillar Hierarchy}
% =====================================================================
The TECT 11-pillar structure (Math54) partitions the theoretical scope as:
\begin{itemize}
\item \textbf{Pillar 1}: Mass gap $m^*$ (Brazovskii-regime locked branch).
\item \textbf{Pillar 2}: IR-scale Lorentz invariance (kinematic bound + shell isotropy).
\item \textbf{Pillar 3}: Emergent spin-2 gravity (one-loop via Ward identity).
\item \textbf{Pillar 4}: Electroweak + GUT gauge group emergence and uniqueness.
\item \textbf{Pillar 5}: Chirality from index-theoretic protection (Dirac fermion
  structure).
\item \textbf{Pillar 6}: Generational structure and SM embedding (Higgs+three-family).
\item \textbf{Pillar 7}: Per-generation quantum consistency (Ward + modular arithmetic).
\item \textbf{Pillar 8}: Lorentz interval-certificate transversality.
\item \textbf{Pillar 9}: Equivalence principle from metric perturbation dynamics.
\item \textbf{Pillar 10}: Planck constant $\hbar$ as cosmic phase-transition
  adiabatic invariant.
\item \textbf{Pillar 11}: Cosmological constant $\Lambda$ via four-sector cancellation.
\end{itemize}

Epoch 7 focuses on Pillars 4-7, each requiring independent mathematical closure
at the proof-sketch (T3) or strong-evidence (T4) level or above before
Stage-1 scorecard publication.

% =====================================================================
\section{Mathematical Contributions}
% =====================================================================

\paragraph{Pillar 4: Gauge-group emergence (Math75, Math80, Math80-AddA-D, Math80-AddE)}

Math75 initiates the Pillar 4 closure programme, addressing three sub-questions
(Q1 equivariant cohomology, Q2 RG-flow derivation, Q3 symplectic reduction) via
the order-parameter structure on the BCC defect manifold. The work is PROVED
CONDITIONAL (T6) on the assumption that the Grassmannian quotient $\mathrm{Gr}(2,5)$
with $\mathrm{SU}(5)$ acting via the two-dimensional representation yields the
Standard Model gauge group via stabiliser theorem. Math80 (strategy framework)
outlines the full gauge-emergence research programme, including defect-bundle
topology on $\mathbb{CP}^2$, equivariant-moduli space structure, and Yukawa
coupling emergence from RGE boundary conditions.

Math80-AddA (defect-moduli theorem) establishes that the two-dimensional
$\mathrm{SU}(5)$ representation moduli on Gr(2,5) parametrises the six
$\mathrm{U}(1)$ charges of the SM. The result is PROVED (T7): all six charge
assignments are derivable from equivariant-action considerations without external
input. Math80-AddB (equivariance closure) proves that the $\mathrm{U}(1)_Y$ charge
distribution respects the required equivariance under the full gauge group,
linking topological and representation-theoretic constraints. Math80-AddC (SO(10)
uniqueness) establishes SO(10) as the unique GUT containing the SM gauge group
via the four-flavour Yukawa structure, PROVED CONDITIONAL (T6) on the
single-shell approximation assumption (caveat documented).
Math80-AddE (Yukawa hierarchy) shows that the hierarchy $y_t > y_b > y_\tau$
emerges from RGE scaling of the BCC elastic modulus across the Margnon transition
scale, PROVED (T7).

Verdict on Pillar 4: T6 PROVED CONDITIONAL on defect-bundle transitivity.

\paragraph{Pillar 5: Chirality from index-theoretic protection (Math76)}

Math76 records the then-current Pillar~5 closure programme: the
three-dimensional Dirac-node structure at the BCC-to-condensate phase
boundary carries a topological invariant (the spectral flow of the
associated Clifford operator) that protects the chirality of the
fermionic zero mode. At the time of writing of this epoch, the result
was tagged PROVED (T7) on the basis of Math10--14 (Pillar~5
five-layer mechanism) plus the Math76 boundary-operator
index-theoretic argument. Subsequent canonical bookkeeping records
the chirality emergence at T7 PROVED with the $16$-zero-mode count
HRR-formula-conditional per Math171-AddA / Top-impact paper TI-1.
The three-dimensional Dirac-node argument employs Atiyah--Singer
index methods applied to the boundary operator $D_\partial$ and the
two-valley structure in the order-parameter manifold; the argument is
recorded here as the then-current closure narrative, not as a
re-derivation in the present paper.

Then-current epoch-7 verdict on Pillar~5: T7 PROVED (chirality
emergence). Current canonical tier per
\texttt{Docs/status/TOE-FACT-SHEET.md}: T7 PROVED with the
HRR-conditional zero-mode count tracked separately.

\paragraph{Pillar 6: Generational structure (Math77)}

Math77 addresses the generational structure of the SM within TECT via a
two-step RGE flow: (i) SU(5) GUT scale breakdown to SM via Pati-Salam
intermediate structure; (ii) family-structure emergence from three-component
order-parameter valleys. The result is PROVED CONDITIONAL (T6) on the
Pati-Salam intermediate symmetry assumption (standard in GUT literature but
introduced here as an explicit hypothesis). The Yukawa coupling hierarchy
$y_t \gg y_b \gg y_\tau$ is derived from RGE running, not postulated.

Verdict on Pillar 6: T6 PROVED CONDITIONAL.

\paragraph{Pillar 7: Quantum consistency (Math47, Math48 carryover + modular constraints)}

Pillar 7 closure combines Math47-Math48 Ward identity closure (carried over from
Epoch 5) with explicit modular-arithmetic constraints on the perturbed metric
$h^{\mu\nu}$ ensuring that quantum-state-count calculations remain non-negative
modulo the perturbative order. The constraint encodes the requirement that the
full gravitational backreaction does not produce negative probabilities at any
order. Verdict: T4 STRONG EVIDENCE (proof sketch complete; full analytical
verification multi-turn project, currently partial-automated in the
Math82-H continuum-audit framework).

\paragraph{Stage-1 Scorecard Synthesis (Math78, audit-rolled-back)}

Math78 attempts a global stage-wide tier assignment for all 11 pillars,
concluding that 4 pillars reach T7 (PROVED), 2 reach T6 (PROVED CONDITIONAL),
1 reaches T5 (CLOSED@1-LOOP), 3 reach T4 (STRONG EVIDENCE), and 1 is
open-negatively-defined (Pillar 10, $\hbar$ origin, provisionally unsolved in
classical framework, leading to the quantum-sector programme). The scorecard is
later (post-Epoch 8) audit-rolled back from "Stage-1 complete" to "Stage-1
PARTIAL-ADVANCED" due to discovery of hidden dependencies between pillars (see
Epoch 8-9 narrative), but the tier-assignment methodology itself is sound and
forms the basis for the refined Stage-2 gate programme.

% =====================================================================
\section{Tier Classification and Audit Trail}
% =====================================================================
This epoch establishes TECT's first canonical tier-assignment framework (see
the internal status-tier nomenclature policy for full tier semantics; binding from 2026-04-29):

\begin{itemize}
\item \textbf{T7 (PROVED)}: Pillar 5 (Math76 chirality), Pillar 3 carryover (gravity),
  Pillars 1 \& 8 carryover (mass gap, transversality).
\item \textbf{T6 (PROVED CONDITIONAL)}: Pillar 4 (Math80-AddA/C, SO(10) uniqueness on
  single-shell), Pillar 6 (Math77, Pati-Salam intermediate stage).
\item \textbf{T5 (CLOSED@1-LOOP)}: Pillar 3 loop-order explicit.
\item \textbf{T4 (STRONG EVIDENCE)}: Pillar 7 (Math47-48 + modular audit), Pillar 9 (equivalence principle),
  Pillar 11 (cosmological constant four-sector).
\item \textbf{T2 (CONJECTURE)}: Pillar 10 ($\hbar$ classical no-go, opens quantum programme).
\end{itemize}

The scorecard assignment is the canonical reference tier-state as of
2026-04-23. Subsequent epochs (8-12) refine this through further audit and
dependent-structure discovery.

% =====================================================================
\section{Discussion and Outlook}
% =====================================================================
Epoch 7 represents the systematic maturation of TECT's proof-tier methodology.
The 11-pillar structure, previously treated as a conceptual checklist, now
becomes a formal tier-assignment matrix, enabling quantitative tracking of
progress toward Stage-1 closure (overall tier = minimum of component tiers).
Stage-1, as defined in Math60, requires all 11 pillars to reach at least T3
(proof sketch); Epoch 7 demonstrates that 8 pillars already exceed T3, with
only Pillar 10 requiring escalation from conjecture to proof-sketch level.

The Pillar 4 closure programme (Math75, Math80) sets the tone for Epochs 8-12:
mathematically deep, making few assumptions, and traceable to numerical evidence
where applicable. The defect-moduli theorem (Math80-AddA) is particularly notable
as the first purely index-theoretic gauge-group derivation in TECT, requiring no
external Higgs sector postulate.

Looking ahead, Epochs 8-9 will pivot toward the quantum sector: $\hbar$-origin
(Pillar 10), quantum-completion gates (GAP-1 through GAP-4), and the
integration of Stage-2 closure conditions (Math60 sub-theorems). The Pillar 4
atlas programme (Epoch 9, Math162-167) will deepen the defect-bundle structure
initiated here, potentially upgrading Pillar 4 from T6 to T7 if the three-patch
Čech-cocycle closure passes audit.

% =====================================================================
\begin{acknowledgments}
The author thanks the Anthropic Claude Opus 4.7 collaborator for systematic
tier-classification discipline and multi-turn audit guidance during the Epoch 7
synthesis. This work is part of TECT, canonical archive at
\url{https://github.com/TECT-OpenPhysics/TECT}.
\end{acknowledgments}

% =====================================================================
\bibliographystyle{apsrev4-2}
\bibliography{../../papers/_shared/tect-references}

\end{document}
